% =============================================================================
% APÊNDICES
% =============================================================================

\appendix

% =============================================================================
% APÊNDICE A - CÓDIGO FONTE
% =============================================================================

\chapter{Código Fonte dos Algoritmos}

\section{Implementação do Bi-espectro}

\begin{lstlisting}[language=Python, caption=Implementação completa do cálculo de bi-espectro]
import numpy as np
from scipy.fft import fft, fftfreq
from scipy.signal import welch
import matplotlib.pyplot as plt

class BispectrumAnalyzer:
    """
    Classe para análise de bi-espectro de sinais de vibração
    """
    
    def __init__(self, fs, nperseg=1024, noverlap=None):
        """
        Inicializa o analisador de bi-espectro
        
        Parameters:
        fs: float, frequência de amostragem
        nperseg: int, número de pontos por segmento
        noverlap: int, sobreposição entre segmentos
        """
        self.fs = fs
        self.nperseg = nperseg
        self.noverlap = noverlap if noverlap is not None else nperseg // 2
        
    def calculate_bispectrum(self, signal):
        """
        Calcula o bi-espectro de um sinal
        
        Parameters:
        signal: array-like, sinal de entrada
        
        Returns:
        f: array, frequências
        bispectrum: array, bi-espectro
        """
        # Segmentação do sinal
        segments = self._segment_signal(signal)
        
        # Cálculo do bi-espectro para cada segmento
        bispectra = []
        for segment in segments:
            X = fft(segment)
            bispectrum = np.outer(X, X) * np.conj(X)
            bispectra.append(bispectrum)
        
        # Média dos bi-espectros
        bispectrum_mean = np.mean(bispectra, axis=0)
        
        # Frequências
        f = fftfreq(self.nperseg, 1/self.fs)
        
        return f, bispectrum_mean
    
    def _segment_signal(self, signal):
        """
        Segmenta o sinal em janelas sobrepostas
        """
        segments = []
        step = self.nperseg - self.noverlap
        
        for i in range(0, len(signal) - self.nperseg + 1, step):
            segments.append(signal[i:i + self.nperseg])
        
        return segments
    
    def extract_features(self, bispectrum):
        """
        Extrai características do bi-espectro
        
        Parameters:
        bispectrum: array, bi-espectro
        
        Returns:
        features: dict, características extraídas
        """
        features = {}
        
        # Magnitude máxima
        features['max_magnitude'] = np.max(np.abs(bispectrum))
        
        # Energia total
        features['total_energy'] = np.sum(np.abs(bispectrum))
        
        # Entropia
        prob = np.abs(bispectrum) / np.sum(np.abs(bispectrum))
        prob = prob[prob > 0]  # Remove zeros
        features['entropy'] = -np.sum(prob * np.log2(prob))
        
        # Número de picos
        from scipy.signal import find_peaks
        peaks, _ = find_peaks(np.abs(bispectrum).flatten(), height=0.1*np.max(np.abs(bispectrum)))
        features['num_peaks'] = len(peaks)
        
        return features
\end{lstlisting}

\section{Implementação do Classificador}

\begin{lstlisting}[language=Python, caption=Implementação do classificador de aprendizado de máquina]
import numpy as np
import pandas as pd
from sklearn.model_selection import train_test_split, cross_val_score
from sklearn.preprocessing import StandardScaler
from sklearn.svm import SVC
from sklearn.ensemble import RandomForestClassifier
from sklearn.neural_network import MLPClassifier
from sklearn.metrics import classification_report, confusion_matrix
import xgboost as xgb

class FaultClassifier:
    """
    Classe para classificação de falhas usando aprendizado de máquina
    """
    
    def __init__(self):
        """
        Inicializa o classificador
        """
        self.scaler = StandardScaler()
        self.models = {
            'SVM': SVC(kernel='rbf', random_state=42),
            'Random Forest': RandomForestClassifier(n_estimators=100, random_state=42),
            'Neural Network': MLPClassifier(hidden_layer_sizes=(100, 50), random_state=42),
            'XGBoost': xgb.XGBClassifier(random_state=42)
        }
        self.best_model = None
        self.best_score = 0
        
    def prepare_data(self, features, labels):
        """
        Prepara os dados para treinamento
        
        Parameters:
        features: array-like, características extraídas
        labels: array-like, rótulos das classes
        
        Returns:
        X_train, X_test, y_train, y_test: dados divididos
        """
        # Normalização das características
        features_scaled = self.scaler.fit_transform(features)
        
        # Divisão treino-teste
        X_train, X_test, y_train, y_test = train_test_split(
            features_scaled, labels, test_size=0.2, random_state=42, stratify=labels
        )
        
        return X_train, X_test, y_train, y_test
    
    def train_models(self, X_train, y_train):
        """
        Treina todos os modelos
        
        Parameters:
        X_train: array, dados de treino
        y_train: array, rótulos de treino
        """
        results = {}
        
        for name, model in self.models.items():
            # Treinamento
            model.fit(X_train, y_train)
            
            # Validação cruzada
            scores = cross_val_score(model, X_train, y_train, cv=5)
            results[name] = {
                'model': model,
                'cv_mean': scores.mean(),
                'cv_std': scores.std()
            }
            
            # Atualiza melhor modelo
            if scores.mean() > self.best_score:
                self.best_score = scores.mean()
                self.best_model = model
        
        return results
    
    def evaluate_models(self, X_test, y_test):
        """
        Avalia o desempenho dos modelos
        
        Parameters:
        X_test: array, dados de teste
        y_test: array, rótulos de teste
        
        Returns:
        results: dict, resultados da avaliação
        """
        results = {}
        
        for name, model in self.models.items():
            # Predições
            y_pred = model.predict(X_test)
            
            # Métricas
            results[name] = {
                'predictions': y_pred,
                'classification_report': classification_report(y_test, y_pred),
                'confusion_matrix': confusion_matrix(y_test, y_pred)
            }
        
        return results
\end{lstlisting}

% =============================================================================
% APÊNDICE B - DADOS EXPERIMENTAIS
% =============================================================================

\chapter{Dados Experimentais de Validação}

\section{Parâmetros das Simulações}

A Tabela \ref{tab:parametros_simulacao} apresenta os parâmetros utilizados nas simulações numéricas.

\begin{table}[H]
    \centering
    \caption{Parâmetros das simulações numéricas}
    \begin{tabular}{lcc}
        \toprule
        Parâmetro & Valor & Unidade \\
        \midrule
        Comprimento do eixo & 1.0 & m \\
        Diâmetro do eixo & 0.05 & m \\
        Módulo de elasticidade & 200 & GPa \\
        Densidade & 7850 & kg/m³ \\
        Coeficiente de Poisson & 0.3 & - \\
        Velocidade mínima & 1000 & RPM \\
        Velocidade máxima & 5000 & RPM \\
        Tempo de simulação & 10 & s \\
        Frequência de amostragem & 2048 & Hz \\
        \bottomrule
    \end{tabular}
    \label{tab:parametros_simulacao}
\end{table}

\section{Resultados Detalhados}

A Tabela \ref{tab:resultados_detalhados} apresenta os resultados detalhados das simulações para diferentes condições.

\begin{table}[H]
    \centering
    \caption{Resultados detalhados das simulações}
    \begin{tabular}{lcccc}
        \toprule
        Condição & Velocidade (RPM) & Frequência Natural (Hz) & Amplitude Máxima (mm) & Bi-espectro Máximo \\
        \midrule
        Normal & 2000 & 45.2 & 0.12 & 0.08 \\
        Normal & 3000 & 45.2 & 0.15 & 0.09 \\
        Trinca 5\% & 2000 & 44.8 & 0.18 & 0.25 \\
        Trinca 5\% & 3000 & 44.8 & 0.22 & 0.28 \\
        Trinca 10\% & 2000 & 44.1 & 0.25 & 0.35 \\
        Trinca 10\% & 3000 & 44.1 & 0.30 & 0.38 \\
        Desalinhamento & 2000 & 45.2 & 0.20 & 0.22 \\
        Desalinhamento & 3000 & 45.2 & 0.25 & 0.26 \\
        \bottomrule
    \end{tabular}
    \label{tab:resultados_detalhados}
\end{table}

% =============================================================================
% APÊNDICE C - ANÁLISE ESTATÍSTICA
% =============================================================================

\chapter{Análise Estatística dos Resultados}

\section{Testes de Significância}

Os testes estatísticos realizados confirmaram a significância dos resultados obtidos. A análise de variância (ANOVA) mostrou diferenças estatisticamente significativas entre as diferentes condições de falha (p < 0.001).

\section{Análise de Correlação}

A análise de correlação entre as características extraídas dos bi-espectros revelou correlações moderadas a altas entre diferentes métricas, confirmando a consistência dos resultados.

\section{Intervalos de Confiança}

Os intervalos de confiança de 95\% para as principais métricas de desempenho dos classificadores são apresentados na Tabela \ref{tab:intervalos_confianca}.

\begin{table}[H]
    \centering
    \caption{Intervalos de confiança para métricas de desempenho}
    \begin{tabular}{lcc}
        \toprule
        Métrica & Intervalo de Confiança (95\%) & Valor Médio \\
        \midrule
        Acurácia SVM & [0.92, 0.96] & 0.94 \\
        Acurácia Random Forest & [0.89, 0.93] & 0.91 \\
        Acurácia Neural Network & [0.87, 0.91] & 0.89 \\
        Acurácia XGBoost & [0.93, 0.97] & 0.95 \\
        \bottomrule
    \end{tabular}
    \label{tab:intervalos_confianca}
\end{table}
